%!TEX root = ../main.tex
%=========================================================

\section{Background}
\label{sec:background}

\michal{Need a good introduction on Kademlia here. Distances/buckets etc. It's required to explain the ticket table, routing etc.}
\etienne{How different is the Ethereum Kademlia from the original published version? I can assume there have been changes in so many years. Is there a reference document for the Ethereum's Kademlia I can read?}
\sergi{From what i remember main difference are num of buckets and use of log distance.  I don't remember any document but Felix will help.}
\sergi{We should also add how kademlia is used for discovery (when lookups are performed, lookup buffer use,  pool of connections, random lookups, etc}

\subsection{Ethereum DHT}
\michal{Borrowed from \href{https://arxiv.org/pdf/1908.10141.pdf}{this paper}}
Ethereum DHT is based on Kademlia - a UDP-based, peer-to-peer distributed hash table (DHT) that is used to locate data stord in a decentralized way~\cite{maymounkov2002kademlia}. Each node is uniquely identified by its randomly generated node ID. A node ID in Ethereum is a marshaled version of node's public key. It is easy to generate many different identities but computationally infeasible to generate a specific one. A data item stored in the DHT is also found by its key, which is simply a hash of the data itself. Node IDs and data keys share the same representation; in the following we use the term ID for both.
Ethereum DHT stores data at nodes whose node ID is the “closest” to the data’s key. Closeness is defined as the bitwise XOR between two IDs, taken as an integer value, i.e.,$d(x, y) = \textit{log}_2 \lfloor x \oplus y \rfloor$.


A DHT node stores its known peers in a routing table divided into buckets which partition the known network based on the local node’s ID (\Cref{fig:kademlia}). Every bucket stores up to $k$ neighbors. Bucket $i$ stores nodes whose distance is in $[2i, 2i+1)$, which effectively corresponds to the length of the common prefix between two node IDs. The routing table stores a detailed view of peers close to the node and a less detailed view of parts of the network located further away.

Items in the DHT and nodes on the network are located by \emph{lookups}. A lookup successively queries nodes that are closest to the desired target ID (key or node ID). Intuitively, a peer closer to the target ID maintains a more detailed view of the target part of the network and provides the initial node with additional peers to query. For instance, the green node in \Cref{fig:kademlia} looking data ID in \emph{bucket 1} will ask known peers in \emph{bucket 1} that will provide additional peers closer to the data ID. Storing data follow the same process.  

\begin{figure}
    \includegraphics[width=0.4\textwidth]{img/kademlia}
    \caption{Ethereum DHT routing table.}
    \label{fig:kademlia}
 \end{figure}

